\section{Sequential receiver and calibration feasibility}
A concrete extension replaces simultaneous separated probes with one active RF chain retuned over 16 blocks from 228 to 316 GHz, each with 16 tones at 1 MHz spacing. A one-sample CP, 30 pilots per 10,000 symbols, both reference and sample, 1 ms settling per hop and whole-frame rounding are charged. The 20 s reference/sample budget contains 19.72 s of transmissions. The WR3.4 converter family supports this frequency conversion range \cite{vdicc}, but the assumed 23 dBm radiated power and 6 dB receiver noise figure remain unverified requirements. Typical intrinsic conversion loss and the quoted small-signal input imply about $-23$ dBm unamplified output before filtering/backoff, leaving at least 46 dB of additional linear gain. This is not demonstrated flight hardware.

Exact Voigt/Mie spectra, refracted paths and emission are recomputed at the selected frequencies. A screened contiguous candidate is rejected after exact joint rank evaluation. With a 0.001 dB correlated differential residual, common gain offset/slope, joint VOC/PM inference and 20 s total, the hopping candidate's response limits are 24.699, 40.176, 4.447 \ugm\ for formaldehyde, methanol and acetonitrile. Independent response draws give approximately 95\% power at each limit and 0.60\% family false alarms. Unknown independent gains at each hop increase acetonitrile's local limit to 548.8 \ugm. Thus retuning requires relative gain calibration, not just a frequency plan.

At 1 \ugm\ acetonitrile, 100 s total gives 3.14\% simulated recall with 0.001 dB residual and 99.99\% with 0.0001 dB residual (10,000 positives per case). At 20 s and 0.0001 dB, recall is 90.20\% (95\% interval 89.60--90.78\%). The 100 s result is a favorable \emph{assumed nonzero calibration} control, not measured stability. More averaging does not overcome the persistent residual. The local requirement for 95\% power at 100 s is approximately 0.000206 dB under the specified correlation model. Arbitrary systematic drift requires a separate protected-threshold budget.

Raw moving controls now include IFFT/CP generation, integer timing acquisition, CP frequency correction, existing-pilot phase tracking, Hamming (7,4) coding and CRC16. The tested initial Doppler prediction error is 100 kHz; a 700 kHz control fails acquisition and is retained. Three bursts per hop sample changing geometry; only generated symbols enter the sensing variance. Across two reference/sample controls, 66,240 packets decode correctly, and passive sensing preserves decoded output. These two trajectories do not estimate recall. Hopping/frame scheduling nevertheless costs 1.382\% relative to a fixed-band modem with the same code and bandwidth. Processing buffers a 10.625 ms frame. Narrowband Doppler per hop and matched reference weather remain assumed; wideband time dilation, fractional timing, oscillator noise and real-time hardware latency remain unvalidated. PM, absolute baseline truth and field compliance are unresolved.
